\writeSectionFrame{\presentationTitle}{Das Problem}
\begin{frame}{Untersysteme}
	\begin{center}
 		\includegraphics[scale=0.2]{\BILDERPFAD/problem1.png}
 	\end{center}
\end{frame}
\note{
	Wir haben irgendeine Art von Untersystem, das aus vielen Klassen besteht. Nun haben wir eine Klasse,
	die auf viele Klassen dieses Untersystems zugreift. Man muss Objekte von all diesen Klassen
	initialisieren, Abhängigkeiten im Auge behalten, Methoden in der korrekten Reihenfolge ausführen usw.
}

\begin{frame}{Fremdbibliotheken}
	\begin{center}
 		\includegraphics[scale=0.3]{\BILDERPFAD/problem2.png}
 	\end{center}
\end{frame}
\note{
	Wir können auf eine Fremdbibliothek zugreifen und diese möchten wir eventuell irgendwann auch
	wieder austauschen. Wenn es aber jetzt viele Abhängigkeiten zu einzelnen Klassen gibt, ist
	ein Austausch nahezu unmöglich.
}

\begin{frame}{Firmenübernahme}
	\begin{center}
 		\includegraphics[scale=0.24]{\BILDERPFAD/problem3.png}
 	\end{center}
\end{frame}
\note{
	Oder ein anderes Szenario. Unsere Firma Schleckermann übernimmt die Konkurrenzfirma
	BestellFix. Beide Unternehmen haben eine Webseite und eine dahinter liegende Backend-Software.
	Wir möchten die Software von BestellFix aus organisatorischen Gründen erst einmal beibehalten,
	weil die Umstellung der über 3000 Filialen einige Jahre dauern wird, aber für den Kunde soll
	es nur noch Schleckermann geben, d.h. es soll auch nur noch eine Webseite geben. Aber unsere
	Webanwendung ist absolut abhängig von unserem eigenen System.
}